\begin{notechapter}
\chapter{Contrastive Learning and the InfoNCE Loss}\label{ch:contrastive-learning}
\chaptersubtitle{Learning representations by recognizing the positive among negatives}

\begin{chapterintro}
Contrastive learning trains an encoder without labels: two views of the same object, such as
two random crops of one image, should get similar representations, and views of different
objects dissimilar ones. The loss that implements this, InfoNCE, is a softmax classifier that
must pick the matching view among $N$ candidates. This chapter derives it from first principles:
the contrastive structure already present in maximum likelihood, the Bayes-optimal
classifier, which is a density ratio, and the lower bound that InfoNCE gives on the mutual
information between the views, with its $\log N$ ceiling. It then analyses the temperature and the
passage from the population loss to the in-batch loss used in practice. You will learn
\begin{itemize}[nosep]
  \item the contrastive setting: encoders, positive pairs, negatives and similarity;
  \item why the gradient of any softmax likelihood pulls the observed item up and pushes the
        model's average down;
  \item InfoNCE as $N$-way classification, its optimal critic $\log\frac{p(y\mid x)}{p(y)} + c(x)$, and the bound $\mathrm I(X; Y)\geq\log N - \mathcal L$;
  \item what the temperature does, and why it is needed with cosine similarities;
  \item in-batch negatives, false negatives, and the related triplet, $N$-pair and supervised
        contrastive losses.
\end{itemize}
\prerequisites{NCE and negative sampling (\cref{ch:nce-negative-sampling}); entropy and the
KL divergence (\cref{ch:cross-entropy}); nested Monte Carlo estimators (\cref{ch:monte-carlo}).}
\end{chapterintro}

\section{The contrastive setting}\label{contrastive-learning:sec:setting}

An encoder $f_\theta:\mathcal X\to\R^d$ maps an input (an image, a sentence, a sound) to a vector, its \emph{representation}.
Without labels, the only supervision available is structure in the data itself, and
contrastive learning uses one kind of structure: pairs of inputs known to share their
meaning. These \emph{positive pairs} are obtained by
\begin{itemize}[nosep]
  \item \emph{data augmentation}: two random transformations (crop, colour jitter, blur) of the same
        image;
  \item \emph{temporal or spatial proximity}: neighbouring frames of a video, nearby words of a
        text (the skip-gram pairs of \cref{ch:word-embeddings});
  \item \emph{multiple modalities}: an image and its caption.
\end{itemize}
Pairs formed from unrelated inputs are \emph{negatives}. A similarity $s(z, z')$ between representations,
usually the cosine $z^\top z'/(\norm z\,\norm{z'})$, is trained to be high on positive pairs and low on negative ones.
Representations are then judged on \emph{downstream} tasks, typically by the accuracy of a
linear classifier trained on top of the frozen encoder.

\begin{pitfall}[Why negatives are needed]
Making positive pairs similar is not enough: the constant encoder $f_\theta(x) = z_0$ makes every pair
perfectly similar and carries no information. Something must push different inputs apart.
Contrastive methods use negatives; non-contrastive methods (\cref{contrastive-practice:sec:collapse})
use architectural asymmetries or explicit variance constraints instead.
\end{pitfall}

\section{The contrastive structure of maximum likelihood}\label{contrastive-learning:sec:mle}

Contrast is not specific to representation learning: it is already present in the gradient of
any softmax model. Let $p_\theta(x) = e^{s_\theta(x)}/Z_\theta$ over a finite set $\mathcal V$, with score $s_\theta$ (an energy-based model with energy $-s_\theta$) and
$Z_\theta = \sum_{x'}e^{s_\theta(x')}$.

\begin{proposition}[Gradient of a softmax log-likelihood]\label{contrastive-learning:prop:mle}
\[
  \nabla_\theta\log p_\theta(x) = g_x - \bar g,\qquad g_x = \nabla_\theta s_\theta(x),\qquad\bar g = \E_{x'\sim p_\theta}[g_{x'}] = \sum_{x'}p_\theta(x')\,g_{x'} .
\]
\end{proposition}

\begin{proof}
$\log p_\theta(x) = s_\theta(x) - \log Z_\theta$ and $\nabla_\theta\log Z_\theta = \frac1{Z_\theta}\sum_{x'}e^{s_\theta(x')}\nabla_\theta s_\theta(x') = \bar g$.
\end{proof}

The first term (the \emph{positive phase}) raises the score of the observed item, the second (the
\emph{negative phase}) lowers the scores of everything in proportion to its current probability. To
see what one gradient-ascent step $\theta\leftarrow\theta + \eta(g_x - \bar g)$ does to the probabilities, expand to first order:
\[
  \Delta\log p_\theta(x')\approx\eta\,(g_{x'} - \bar g)^\top(g_x - \bar g)\qquad\text{for every }x' .
\]
Three consequences follow.
\begin{itemize}
  \item The observed item gains: $\Delta\log p_\theta(x)\approx\eta\norm{g_x - \bar g}^2\geq0$.
  \item An item $x'$ gains when its centred gradient $g_{x'} - \bar g$ points in the same direction as that of $x$, and
        loses when it points the opposite way: ``similar'' items in the sense of the
        parameterization rise together.
  \item Probability is conserved: $\sum_{x'}p_\theta(x')\,\Delta\log p_\theta(x')\approx\eta\bigl(\sum_{x'}p_\theta(x')(g_{x'} - \bar g)\bigr)^\top(g_x - \bar g) = 0$, which is the first-order
        form of $\sum_{x'}\Delta p_\theta(x') = 0$.
\end{itemize}
Note that it is $\log p$, not the score $s$, that behaves this way: the score change $\eta\,g_{x'}^\top(g_x - \bar g)$ can be
negative even for the observed item, since scores are only defined up to a common constant.

\paragraph{The cost of the negative phase.} The average $\bar g$ is a sum over $\mathcal V$, which is the whole difficulty
of training softmax models over large sets (\cref{ch:hierarchical-softmax}). NCE and negative
sampling (\cref{ch:nce-negative-sampling}) replace it with a few sampled noise items: NCE recovers
normalized log-probabilities, negative sampling a shifted pointwise mutual information. InfoNCE,
next, does the same with a softmax over sampled candidates, and estimates the density ratio
$p(x, y)/(p(x)p(y))$ between paired variables.

\section{InfoNCE}\label{contrastive-learning:sec:infonce}

\paragraph{Setting.} Let $(X, Y)$ be a pair of random variables with joint density $p(x, y)$, for instance two views of
the same image, and let $r(x, y) = \frac{p(x, y)}{p(x)p(y)} = \frac{p(y\mid x)}{p(y)}$ be their \emph{density ratio}. Its logarithm is the pointwise
mutual information, and its average under the joint distribution is the mutual information,
\[
  \mathrm I(X; Y) = \E_{p(x,y)}\bigl[\log r(X, Y)\bigr] = \mathrm{KL}\bigl(p(x, y)\,\Vert\,p(x)p(y)\bigr) .
\]
Draw $x\sim p(x)$, a positive $y_0\sim p(y\mid x)$ and $N - 1$ negatives $y_1,\dots,y_{N-1}\sim p(y)$, independently, and place the positive at a
uniformly random position among the $N$ candidates. Which candidate is the positive?

\begin{proposition}[Bayes-optimal identification]\label{contrastive-learning:prop:bayes}
Given $x$ and the candidates, the probability that candidate $j$ is the positive is
\[
  \Prob(\text{positive} = j\mid x, y_0,\dots,y_{N-1}) = \frac{r(x, y_j)}{\sum_{k=0}^{N-1}r(x, y_k)} .
\]
\end{proposition}

\begin{proof}
If candidate $j$ is the positive, the candidates have density $p(y_j\mid x)\prod_{k\neq j}p(y_k)$. With a uniform prior over $j$,
Bayes' rule gives a posterior proportional to this density. Dividing by $\prod_kp(y_k)$, which does not depend on $j$,
leaves $p(y_j\mid x)/p(y_j) = r(x, y_j)$.
\end{proof}

The ratio $r(x, y_j)$ is computed at the sampled point and is not $1$ for a negative: independence of $x$ and $y_j$ as
random variables does not make the ratio equal to $1$ at the points drawn, only its expectation under
$p(x)p(y)$ (\cref{contrastive-learning:exr:1}).

\paragraph{The loss.} A \emph{critic} $f(x, y)$, for example $f(x, y) = \cos(f_\theta(x), f_\theta(y))/\tau$, defines the classifier $\softmax_j f(x, y_j)$. Its
cross-entropy, with the positive at index $0$, is the \emph{InfoNCE} loss (van den Oord, Li and Vinyals,
\emph{Representation Learning with Contrastive Predictive Coding}, 2018):
\[
  \mathcal L_N(f) = \E\Bigl[-\log\frac{e^{f(x, y_0)}}{\sum_{j=0}^{N-1}e^{f(x, y_j)}}\Bigr],
\]
the expectation being over $x$, $y_0\sim p(y\mid x)$ and the negatives. It is the softmax cross-entropy of
\cref{bigram-models:sec:gradient} over the $N$ candidates.

\begin{theorem}[The optimal critic]\label{contrastive-learning:thm:critic}
Over all critics, $\mathcal L_N$ is minimized by $f^*(x, y) = \log r(x, y) + c(x)$, for any function $c$ of $x$ alone, and (for $N\geq2$) only by
critics of this form, up to sets of probability zero.
\end{theorem}

\begin{proof}
The loss is the expected cross-entropy between the Bayes posterior of
\cref{contrastive-learning:prop:bayes} and the classifier $\softmax_jf(x, y_j)$ (by symmetry, the position of the positive
does not change the loss). For each $(x, y_{0:N-1})$, the cross-entropy is minimized exactly when the two
distributions over $j$ are equal (Gibbs' inequality, \cref{ch:cross-entropy}), that is when $f(x, y_j) - \log r(x, y_j)$ does not
depend on $j$. Since the candidates are arbitrary draws, this forces $f(x, y) - \log r(x, y)$ to depend on $x$ only.
\end{proof}

A critic of the form $\cos(f_\theta(x), f_\theta(y))/\tau$ can only approximate $f^*$, but the theorem says what training aims at: the
similarity between two representations should be the pointwise mutual information of the
inputs, up to a per-anchor constant.

\begin{theorem}[InfoNCE bounds the mutual information]\label{contrastive-learning:thm:bound}
For every critic $f$,
\[
  \mathrm I(X; Y)\geq\log N - \mathcal L_N(f) .
\]
\end{theorem}

\begin{proof}
\emph{Step 1, a variational bound.} For distributions $P$, $Q$ and any function $T$,
$\mathrm{KL}(P\Vert Q)\geq\E_P[T] - \E_Q[e^{T-1}]$ (the bound of Nguyen, Wainwright and Jordan, 2010, from the convex
duality of the KL divergence; \cref{contrastive-learning:exr:2} proves it).

\emph{Step 2, apply it to the whole batch.} Let $\mathbf Y = (y_0,\dots,y_{N-1})$, with $(x, y_0)\sim p(x, y)$ and the negatives
independent. Since the negatives are independent of $(x, y_0)$, $\mathrm I(X;\mathbf Y) = \mathrm I(X; Y)$. Take $P$ the joint law of $(x, \mathbf Y)$, $Q$
the product of its marginals, and
\[
  T(x,\mathbf Y) = 1 + \log\frac{e^{f(x, y_0)}}{m(x,\mathbf Y)},\qquad m(x,\mathbf Y) = \frac1N\sum_{j=0}^{N-1}e^{f(x, y_j)} .
\]
Under $Q$ the $y_j$ are i.i.d.\ from $p(y)$ and independent of $x$, so the $N$ expectations $\E_Q\bigl[e^{f(x, y_k)}/m\bigr]$, $k = 0,\dots,N-1$,
are equal; their average is $\E_Q[m/m] = 1$, hence $\E_Q[e^{T-1}] = \E_Q\bigl[e^{f(x, y_0)}/m\bigr] = 1$.

\emph{Step 3, conclude.} $\mathrm I(X; Y)\geq\E_P[T] - 1 = \E\bigl[\log\frac{e^{f(x, y_0)}}{\frac1N\sum_je^{f(x, y_j)}}\bigr] = \log N - \mathcal L_N(f)$.
\end{proof}

\begin{pitfall}[A tempting but wrong proof]
A frequently reproduced argument plugs in the optimal critic, claims that $r(x, y_j) = 1$ for the
negatives ``because $y_j$ is independent of $x$'', and uses $\E_{p(x,y)}[r] = 1$. Both claims are false: $r(x, y_j)$ is a
random variable with mean $1$ under $p(x)p(y)$, and under the joint distribution
$\E_{p(x,y)}[r] = 1 + \chi^2\bigl(p(x, y)\Vert p(x)p(y)\bigr)\geq1$. Moreover, a statement about the optimal critic could not give a bound valid for
every critic, which is what training needs. The proof above uses neither claim.
\end{pitfall}

\paragraph{The $\log N$ ceiling.} Since $\mathcal L_N\geq0$, the bound never exceeds $\log N$: with $N$ candidates, InfoNCE can certify at
most $\log N$ nats of mutual information, however strongly $X$ and $Y$ depend on each other. McAllester
and Stratos (\emph{Formal Limitations on the Measurement of Mutual Information}, 2020) show that
this is not a weakness of InfoNCE alone: any distribution-free high-confidence lower bound on
mutual information computed from $N$ samples is at most $O(\log N)$.

\begin{example}[Correlated Gaussians]\label{contrastive-learning:ex:gauss}
Let $(X, Y)$ be standard normal with correlation $\rho$, so that $\mathrm I(X; Y) = -\frac12\log(1 - \rho^2)$ and $\log r$ is an explicit
quadratic function. With the optimal critic $f^* = \log r$, Monte Carlo estimates of $\log N - \mathcal L_N(f^*)$ (\cref{contrastive-learning:fig:gauss})
show both regimes. For $\rho = 0.9$, $\mathrm I = 0.830$ nats, and the bound rises from $0.337$ with $N = 2$ to $0.809$ with $N = 64$: a
moderate mutual information is reached with tens of negatives. For $\rho = 0.999$, $\mathrm I = 3.108$ nats, and the bound
stays close to the ceiling $\log N$ for small $N$ ($0.654$ for $N = 2$, where $\log2 = 0.693$; $1.831$ for $N = 8$, where $\log8 = 2.079$),
and still falls short with $N = 1024$ ($3.087$). The optimal critic cannot do better: the gap is a
property of the bound.
\end{example}

\begin{figure}[ht]
\centering
\begin{tikzpicture}
\begin{axis}[width=0.75\textwidth, height=5.8cm, xmode=log, log basis x=2, xlabel={number of candidates $N$}, ylabel={nats},
  xmin=2, xmax=1024, ymin=0, ymax=3.6, grid=major, grid style={gray!20}, legend style={font=\scriptsize, at={(0.98,0.5)}, anchor=east},
  tick label style={font=\scriptsize}, label style={font=\small}]
\addplot[thick, gray, dashed, domain=2:1024, samples=50] {ln(x)};
\addlegendentry{$\log N$}
\addplot[thick, bookpitfall, mark=*, mark size=1.2pt] coordinates {(2,0.654) (4,1.274) (8,1.831) (16,2.291) (32,2.626) (64,2.839) (128,2.958) (256,3.025) (512,3.061) (1024,3.087)};
\addlegendentry{bound, $\rho = 0.999$}
\addplot[thin, bookpitfall, densely dotted, domain=2:1024, forget plot] {3.108};
\addplot[thick, bookaccent, mark=*, mark size=1.2pt] coordinates {(2,0.337) (4,0.553) (8,0.679) (16,0.750) (32,0.787) (64,0.809) (128,0.819) (256,0.821) (512,0.824) (1024,0.835)};
\addlegendentry{bound, $\rho = 0.9$}
\addplot[thin, bookaccent, densely dotted, domain=2:1024, forget plot] {0.830};
\end{axis}
\end{tikzpicture}
\caption{The InfoNCE bound $\log N - \mathcal L_N(f^*)$ with the optimal critic for correlated Gaussians, against the true mutual
information (dotted). Monte Carlo estimates with $2\times10^7/N$ draws (at least $2\,000$), standard errors below
$0.007$. For a large mutual information the bound is limited by $\log N$.}
\label{contrastive-learning:fig:gauss}
\end{figure}

\begin{remark}[Mutual information is not the whole story]
Maximizing a mutual-information bound does not by itself produce good representations:
any invertible function of a representation has the same mutual information with the other
view, yet linear classifiers may perform very differently on the two (Tschannen et al., \emph{On
Mutual Information Maximization for Representation Learning}, 2020). The success of
contrastive learning depends as much on the critic family (normalized dot products), the
encoder architecture and the augmentations as on the bound.
\end{remark}

\section{Temperature}\label{contrastive-learning:sec:temperature}

With cosine similarities, the critic is $f(x, y) = s(x, y)/\tau$ with $s\in[-1, 1]$ and a temperature $\tau > 0$:
\[
  \ell = -\log\frac{e^{s_0/\tau}}{\sum_{j=0}^{N-1}e^{s_j/\tau}},\qquad s_j = s(x, y_j) .
\]

\begin{proposition}[Effect of the temperature]\label{contrastive-learning:prop:temperature}
\begin{enumerate}[nosep]
  \item If $s_0 > s_j$ for all $j\geq1$, then $\ell\to0$ as $\tau\to0^+$; as $\tau\to\infty$, $\ell\to\log N$ whatever the similarities.
  \item With $s\in[-1, 1]$, $\ell\geq\log\bigl(1 + (N - 1)e^{-2/\tau}\bigr) > 0$: with $\tau = 1$ and $N = 4096$, the loss cannot go below $6.32$.
  \item $\nabla_\theta\ell = \frac1\tau\bigl(\sum_jp_j\nabla_\theta s_j - \nabla_\theta s_0\bigr)$ with $p_j = \softmax_j(s_j/\tau)$: each negative is pushed with weight $p_j$.
\end{enumerate}
\end{proposition}

\begin{proof}
(1) $\ell = \log\bigl(1 + \sum_{j\geq1}e^{(s_j - s_0)/\tau}\bigr)$; each exponent tends to $-\infty$ as $\tau\to0^+$ and to $0$ as $\tau\to\infty$. (2) Each term
$e^{(s_j - s_0)/\tau}\geq e^{-2/\tau}$ since $s_j - s_0\geq-2$. For $\tau = 1$, $N = 4096$: $\log(1 + 4095e^{-2}) = 6.32$. (3) As in
\cref{contrastive-learning:prop:mle}, with $f_j = s_j/\tau$.
\end{proof}

Point (2) explains why a small temperature is needed with cosine similarities: the
logits are confined to an interval of width $2/\tau$, and with $\tau = 1$ the classifier cannot become confident.
Point (3) explains its other effect: as $\tau$ decreases the weights $p_j$ concentrate on the negatives with the
highest similarity, the \emph{hard negatives}, which dominate the gradient; as $\tau$ increases all negatives
are pushed alike. Typical values are $\tau = 0.07$ (memory-bank methods and MoCo), $0.1$ to $0.5$ (SimCLR), or a
learned temperature (CLIP), see \cref{ch:contrastive-practice}.

\section{From the population loss to a batch}\label{contrastive-learning:sec:batch}

In practice the candidates come from the mini-batch. Given $B$ positive pairs $(x_i, y_i)$, the other $y_j$ of the
batch serve as negatives for $x_i$:
\[
  \hat{\mathcal L} = -\frac1B\sum_{i=1}^B\log\frac{e^{f(x_i, y_i)}}{\sum_{j=1}^Be^{f(x_i, y_j)}} .
\]
\begin{itemize}
  \item \emph{What it estimates.} If the $B$ pairs are i.i.d., each term has exactly the distribution of the InfoNCE
        loss with $N = B$ candidates, so $\E[\hat{\mathcal L}] = \mathcal L_B$: the batch loss is an unbiased estimate of the $N = B$ objective. (The
        terms are dependent, since they share candidates, which affects the variance, not the
        mean.) It is \emph{not} an unbiased estimate of the objective with more candidates: $\mathcal L_N$ is the
        expectation of a nested quantity, a logarithm of a sum, and changes with $N$, as in the nested
        estimators of \cref{monte-carlo:sec:nested}.
  \item \emph{Symmetric version.} When both sides are views of the same kind, the loss is averaged with
        the same loss with the roles of $x$ and $y$ exchanged (SimCLR, CLIP).
  \item \emph{False negatives.} The negatives are drawn from $p(y)$, not from ``inputs with a different
        meaning'': a negative may show an object of the same class as the anchor. With $C$
        equally frequent classes, a batch of $B$ images gives each anchor about $(B - 1)/C$ same-class
        images among its negatives: for ImageNet ($C = 1000$) and $B = 4096$, about $4$ images, $8$ views in the
        two-view setting of SimCLR. They are pushed away although they should not be. Debiased
        variants correct the loss for this effect (Chuang et al., \emph{Debiased Contrastive
        Learning}, 2020).
\end{itemize}

\section{Related losses}\label{contrastive-learning:sec:related}

\paragraph{Triplet loss.} With an anchor $a$, a positive $p$, a negative $n$, a distance $d$ and a margin $m > 0$, the
triplet loss is $\max\bigl(0,\ d(a, p) - d(a, n) + m\bigr)$ (Schroff, Kalenichenko and Philbin, \emph{FaceNet}, 2015). InfoNCE with one
negative and the critic $f = -d^2/\tau$ is $\log\bigl(1 + e^{(d(a,p)^2 - d(a,n)^2)/\tau}\bigr)$, a smooth version of the triplet loss with squared distances
and no margin (\cref{contrastive-learning:exr:4}).

\paragraph{$N$-pair loss.} $\log\bigl(1 + \sum_{j=1}^{N-1}e^{f(a, n_j) - f(a, p)}\bigr)$ (Sohn, 2016) is exactly InfoNCE with $N - 1$ negatives, written with the
positive's logit subtracted.

\paragraph{Supervised contrastive loss.} When labels are available, every other sample of the same
class is a positive (Khosla et al., \emph{Supervised Contrastive Learning}, 2020): for anchor $i$ with
positives $P(i)$,
\[
  \ell_i = -\frac1{\abs{P(i)}}\sum_{p\in P(i)}\log\frac{e^{f(x_i, x_p)}}{\sum_{a\neq i}e^{f(x_i, x_a)}},
\]
which removes the false negatives by construction.

\begin{chaptersummary}
\begin{itemize}
  \item Contrastive learning trains an encoder so that positive pairs (views sharing their meaning)
        are similar and negatives dissimilar; negatives prevent the trivial constant encoder.
  \item Every softmax likelihood has the gradient $g_x - \E_{p_\theta}[g]$: positive phase minus negative phase.
        A step raises $\log p$ of the observed item by $\eta\norm{g_x - \bar g}^2$ and conserves probability.
  \item InfoNCE is the cross-entropy of identifying the positive among $N$ candidates; the Bayes
        posterior is $r(x, y_j)/\sum_kr(x, y_k)$ and the optimal critic is $\log r(x, y) + c(x)$, the pointwise mutual information.
  \item $\mathrm I(X; Y)\geq\log N - \mathcal L_N(f)$ for every critic, so the certified mutual information never exceeds $\log N$.
  \item The temperature scales the logits: needed for confident classification with cosine
        similarities, and it controls how much the hardest negatives dominate the gradient.
  \item The in-batch loss is unbiased for the $N = B$ objective; it contains false negatives.
\end{itemize}
\end{chaptersummary}

\section*{Exercises}
\addcontentsline{toc}{section}{Exercises}

\begin{exercise}\label{contrastive-learning:exr:1}
Show that $\E_{p(x)p(y)}[r(X, Y)] = 1$ and $\E_{p(x,y)}[r(X, Y)] = 1 + \chi^2\bigl(p(x, y)\Vert p(x)p(y)\bigr)$, where $\chi^2(P\Vert Q) = \int(p - q)^2/q$. For the
Gaussian pair of \cref{contrastive-learning:ex:gauss} with $\rho = 0.5$, is $r(x, y) = 1$ at $x = y = 1$?
\solutionlink{contrastive-learning:exr:1}
\end{exercise}

\begin{exercise}\label{contrastive-learning:exr:2}
Prove the bound $\mathrm{KL}(P\Vert Q)\geq\E_P[T] - \E_Q[e^{T-1}]$ for discrete $P, Q$, with equality for $T = 1 + \log(p/q)$.
\hint{Use $\log u\geq1 - 1/u$, or maximize $t\mapsto p\,t - q\,e^{t-1}$ pointwise.}
\solutionlink{contrastive-learning:exr:2}
\end{exercise}

\begin{exercise}\label{contrastive-learning:exr:3}
Show that $\mathcal L_N(f)\leq\log N$ when $f$ is constant, and that the bound of \cref{contrastive-learning:thm:bound} is then $0$. What does a
loss above $\log N$ mean?
\solutionlink{contrastive-learning:exr:3}
\end{exercise}

\begin{exercise}\label{contrastive-learning:exr:4}
Derive the expression of InfoNCE with one negative and $f = -d^2/\tau$, and compare its behaviour with the
triplet loss when $d(a, n)^2 - d(a, p)^2$ is large and when it is negative.
\solutionlink{contrastive-learning:exr:4}
\end{exercise}

\begin{exercise}\label{contrastive-learning:exr:5}
In \cref{contrastive-learning:prop:temperature}, compute the minimal loss for $\tau = 0.1$ and $N = 4096$. For a batch with logits
$s/\tau$, $s = (0.9, 0.8, 0.1, -0.5)$ (positive first), compute the weights $p_j$ for $\tau = 1$ and $\tau = 0.1$, and say which negative
dominates the gradient.
\solutionlink{contrastive-learning:exr:5}
\end{exercise}
\end{notechapter}
